\providecommand{\todo}[1]{\textcolor{red}{\textbf{[TODO: #1]}}}
\clearpage
\phantomsection
\section*{Chương 5. Trừ trung bình cục bộ, với DC-blocker tùy chọn}
\addcontentsline{toc}{section}{Chương 5. Trừ trung bình cục bộ, với DC-blocker tùy chọn}

\begin{learningobj}
\begin{itemize}[leftmargin=1.5em, itemsep=2pt]
    \item Thấu hiểu bản chất vật lý và nguồn gốc toán học của thành phần một chiều (DC offset) trong tín hiệu âm thanh; chứng minh chặt chẽ cơ chế độ lệch $+c^2$ nâng bổng đường tự tương quan chuẩn hóa và làm sụp đổ ngưỡng quyết định hữu thanh (\texttt{VOICING\_THRESHOLD} $= 0{,}45$).
    \item Khám phá mã nguồn Praat 7.0.03 (\texttt{Sound\_to\_Pitch.cpp}, dòng 59--66), phân tích vì sao Boersma tính trung bình cục bộ trên đúng 426 mẫu ($2 \times 213$ mẫu quanh tâm khung) thay vì 640 hay 638 mẫu; chứng minh tính tối ưu khi đọc trực tiếp từ bộ đệm vòng (Ring Buffer) mà không tốn thêm bộ nhớ RAM.
    \item Nắm vững kỹ thuật chia cho hằng số bằng phép nhân số nguyên nghịch đảo; vạch rõ sai lầm của Argon ($M = 26.214, k = 24$ gây lỗi $25\%$), dẫn xuất thuật toán Golden v2 ($M = 10.082.083, k = 32$) kết hợp độ lệch không âm $B = 426 \times 2^{15}$, bảo đảm khớp bit tuyệt đối $100\%$ trên toàn bộ $27.917.911$ giá trị tổng vật lý.
    \item Khảo sát chuyên sâu khối IP Intel \texttt{LPM\_MULT} theo tài liệu \emph{Integer Arithmetic IP Cores User Guide} (UG-01069) và kiến trúc khối DSP biến đổi độ chính xác trong \emph{Cyclone V Device Handbook} (CV-52003); giải thích tường tận vì sao phép nhân $25 \times 24$-bit không dấu vừa vặn trong đúng 1 khối DSP 27$\times$27.
    \item Hiện thực hóa module phần cứng Verilog-2001 (\texttt{ch05\_local\_mean.v}) và bộ lọc DC-blocker tùy chọn bậc 1 (\texttt{ch05\_dc\_blocker.v}, $\alpha = 1 - 2^{-8}$); giải mã cấu trúc thanh ghi tích lũy 36-bit chống hiện tượng tràn số và chu trình giới hạn (limit cycle).
    \item Thực thi kiểm chứng bit-đối-bit và chu-kỳ-đối-chu-kỳ giữa mô hình song sinh Python và mô phỏng HDL trên 205 khung tín hiệu thực nghiệm và các trường hợp góc biên khắc nghiệt.
\end{itemize}
\end{learningobj}

\subsection*{5.1. Bản chất và tác hại của mức DC trong tự tương quan}
\addcontentsline{toc}{subsection}{5.1. Bản chất và tác hại của mức DC trong tự tương quan}

\subsubsection*{Vấn đề thực tế: Sự thật phũ phàng về độ lệch một chiều và tiếng thì thầm bị biến thành tiếng nói}
Giả sử một hệ thống thu âm giọng nói có micro hoặc mạch khuếch đại tiền khuếch đại (preamplifier) bị lệch điện áp phân cực nhẹ. Dòng âm thanh số hóa $16\text{ bit}$ thu được không dao động quanh mức $0$ mà bị nâng lên một mức một chiều $c = +1000$ đơn vị biên độ (chỉ chiếm khoảng $3\%$ dải toàn thang $32.768$, tai người hoàn toàn không nghe thấy tiếng ù này). 

Khi người nói dừng phát âm để chuyển sang một đoạn thì thầm hoặc một khoảng lặng vô thanh (vùng unvoiced hoặc silence), năng lượng tín hiệu thực chất chỉ là nhiễu nền ngẫu nhiên với độ lệch chuẩn $\sigma = 1000$. Thuật toán tự tương quan trên FPGA được kích hoạt để phân tích khung $638\text{ mẫu}$. Nếu mạch phần cứng bỏ qua bước khử trung bình và đưa trực tiếp tín hiệu có mức DC vào khối nhân cửa sổ Hanning rồi tính tự tương quan $R(\tau)$, điều gì sẽ xảy ra?

Thay vì nhận diện đây là khoảng lặng với $r(\tau) \approx 0$ và phân loại đúng là vô thanh (Unvoiced), giá trị tự tương quan chuẩn hóa tại mọi độ trễ $\tau$ bị kéo vọt lên trên $0{,}60$. Con số này vượt xa ngưỡng hữu thanh chuẩn của Praat (\texttt{VOICING\_THRESHOLD} $= 0{,}45$). Hệ thống kết luận sai lầm rằng đây là một nguyên âm hữu thanh (Voiced) với độ tin cậy rất cao, sinh ra một cao độ $F_0$ giả tạo và làm bùng nổ tỷ lệ lỗi quyết định hữu thanh (VDE). Toàn bộ thuật toán sụp đổ chỉ vì một mức điện áp tĩnh không nhìn thấy bằng mắt thường.

\subsubsection*{Dẫn xuất từ nguyên lý gốc: Cơ chế xuất hiện thành phần $+c^2$ nâng bổng đường tự tương quan}
Để hiểu vì sao mức DC tĩnh $c$ lại gây ra hậu quả nghiêm trọng đến vậy, ta phân tích biểu thức toán học từ nguyên lý đầu tiên. Giả sử đoạn tín hiệu âm thanh thực $x[n]$ bao gồm thành phần sóng âm dao động thực thụ $s[n]$ (có giá trị kỳ vọng bằng $0$) cộng với một mức một chiều không đổi $c$:
\begin{equation}
x[n] = s[n] + c, \quad \text{với } \frac{1}{N} \sum_{n=0}^{N-1} s[n] \approx 0
\end{equation}

Khi áp dụng hàm cửa sổ thời gian $w[n]$ (cửa sổ Hanning độ dài $N = 638$), tín hiệu sau cửa sổ hóa là:
\begin{equation}
x_w[n] = x[n] \cdot w[n] = (s[n] + c) \cdot w[n] = s_w[n] + c \cdot w[n]
\end{equation}
trong đó $s_w[n] = s[n] w[n]$ là tín hiệu sóng âm đã nhân cửa sổ. Hàm tự tương quan ngắn hạn chưa chuẩn hóa tại độ trễ $\tau$ được tính bằng tích vô hướng trượt:
\begin{equation}
R_{xw}(\tau) = \sum_{n=0}^{N-1-\tau} x_w[n] \cdot x_w[n+\tau] = \sum_{n=0}^{N-1-\tau} \Big[ s_w[n] + c \cdot w[n] \Big] \Big[ s_w[n+\tau] + c \cdot w[n+\tau] \Big]
\end{equation}

Khai triển biểu thức tích dưới dấu tổng, ta thu được bốn số hạng riêng biệt:
\begin{equation}
R_{xw}(\tau) = \underbrace{\sum_{n=0}^{N-1-\tau} s_w[n] s_w[n+\tau]}_{R_{sw}(\tau)} + c \sum_{n=0}^{N-1-\tau} s_w[n] w[n+\tau] + c \sum_{n=0}^{N-1-\tau} s_w[n+\tau] w[n] + c^2 \underbrace{\sum_{n=0}^{N-1-\tau} w[n] w[n+\tau]}_{R_w(\tau)}
\end{equation}
trong đó $R_{sw}(\tau)$ là tự tương quan của tín hiệu sạch, còn $R_w(\tau)$ là tự tương quan năng lượng của chính hàm cửa sổ Hanning.

Vì $s[n]$ là tín hiệu dao động xoay chiều có giá trị trung bình bằng $0$ và không tương quan với mức tĩnh $c$, các số hạng chéo tích giữa $s_w$ và $w$ sẽ triệt tiêu lẫn nhau khi lấy tổng trên số lượng mẫu đủ lớn:
\begin{equation}
c \sum_{n=0}^{N-1-\tau} s_w[n] w[n+\tau] \approx 0, \quad c \sum_{n=0}^{N-1-\tau} s_w[n+\tau] w[n] \approx 0
\end{equation}

Khi đó, tự tương quan của tín hiệu có mức DC thu gọn thành:
\begin{equation}
R_{xw}(\tau) \approx R_{sw}(\tau) + c^2 \cdot R_w(\tau)
\end{equation}

Hãy chú ý đến bước chuẩn hóa cửa sổ ở Chương 8: để loại bỏ ảnh hưởng suy giảm năng lượng của hàm cửa sổ Hanning ở các độ trễ lớn, thuật toán Praat bắt buộc phải chia hàm tự tương quan đo được cho hàm tự tương quan của cửa sổ:
\begin{equation}
r_x(\tau) = \frac{R_{xw}(\tau)}{R_w(\tau)} \approx \frac{R_{sw}(\tau) + c^2 \cdot R_w(\tau)}{R_w(\tau)} = \frac{R_{sw}(\tau)}{R_w(\tau)} + c^2 = r_s(\tau) + \mathbf{c^2}
\end{equation}

\begin{physicalinsight}
\textbf{Khám phá then chốt:} Sau khi chuẩn hóa chia cho $R_w(\tau)$, thành phần một chiều $c$ biến thành một \textbf{hằng số dương tuyệt đối} $+c^2$ được cộng phẳng vào mọi độ trễ $\tau$!
\begin{itemize}[leftmargin=1.5em, itemsep=2pt]
    \item Vì $c^2 \ge 0$, mức DC không bao giờ làm giảm mà luôn \textbf{nâng bổng toàn bộ đáy và đỉnh} của đường cong tự tương quan lên cao.
    \item Đối với âm thanh vô thanh hoặc đoạn lặng, $r_s(\tau) \approx 0$. Tuy nhiên, do bị cộng thêm $+c^2$, giá trị tự tương quan chuẩn hóa theo tỷ lệ năng lượng trở thành $r_x(\tau)/r_x(0) \approx c^2 / (\sigma^2 + c^2)$. Khi năng lượng DC $c^2$ tương đương phương sai nhiễu $\sigma^2$, tỷ số này đạt $0{,}50 > 0{,}45$, khiến bộ so sánh voicing khẳng định chắc chắn rằng tín hiệu có tuần hoàn!
\end{itemize}
\end{physicalinsight}

\subsubsection*{Tính toán bằng tay: Ví dụ sóng vuông 8 mẫu cộng mức DC +1000}
Để minh chứng tường minh công thức $r_x(\tau) = r_s(\tau) + c^2$, ta xét một ví dụ số học thu nhỏ gồm $N = 8$ mẫu được tính toán độc lập bằng kịch bản \texttt{verify\_v3/ch05\_dc\_effect.py}. 

Xét tín hiệu sóng vuông đối xứng chu kỳ $4$ mẫu:
\begin{equation}
s[n] = [-1000, -1000, 1000, 1000, -1000, -1000, 1000, 1000]
\end{equation}
Tín hiệu này có giá trị trung bình thực sự bằng $0$ ($\sum s[n] = 0$). Ta cộng thêm mức DC $c = +1000$, tạo thành tín hiệu thực tế:
\begin{equation}
x[n] = s[n] + 1000 = [0, 0, 2000, 2000, 0, 0, 2000, 2000]
\end{equation}

Áp dụng hai trường hợp cửa sổ: cửa sổ chữ nhật ($w[n] = 1$) và cửa sổ Hanning chuẩn Praat ($w[i] = 0{,}5 - 0{,}5\cos(2\pi i / 9)$, $i = 1 \dots 8$). Bảng~\ref{tab:ch05_dc_hand_calc} trình bày chi tiết kết quả tính toán tự tương quan và tự tương quan chuẩn hóa cho các độ trễ $\tau = 0 \dots 4$.

\begin{table}[H]
\centering
\caption{Tính toán bằng tay tác hại của mức DC $c = +1000$ lên tự tương quan ($N = 8$ mẫu, trích từ log \texttt{verify\_v3/ch05\_dc\_effect.log})}
\label{tab:ch05_dc_hand_calc}
\footnotesize
\begin{tabular}{crrrrrr}
\toprule
\textbf{Độ trễ $\tau$} & \textbf{$R_s(\tau)$} & \textbf{$R_x(\tau)$} & \textbf{$R_w(\tau)$} & \textbf{$r_s(\tau) = R_s/R_w$} & \textbf{$r_x(\tau) = R_x/R_w$} & \textbf{Độ lệch $r_x - r_s$} \\
\midrule
\multicolumn{7}{l}{\textbf{Trường hợp 1: Cửa sổ chữ nhật ($w[n] = 1$)}} \\
0 & $8.000.000{,}0$ & $16.000.000{,}0$ & $8{,}0000$ & $1.000.000{,}00$ & $2.000.000{,}00$ & $\mathbf{+1.000.000{,}00}$ ($= c^2$) \\
1 & $1.000.000{,}0$ & $8.000.000{,}0$  & $7{,}0000$ & $142.857{,}14$   & $1.142.857{,}14$ & $\mathbf{+1.000.000{,}00}$ ($= c^2$) \\
2 & $-6.000.000{,}0$ & $0{,}0$         & $6{,}0000$ & $-1.000.000{,}00$ & $0{,}00$         & $\mathbf{+1.000.000{,}00}$ ($= c^2$) \\
3 & $-1.000.000{,}0$ & $4.000.000{,}0$  & $5{,}0000$ & $-200.000{,}00$  & $800.000{,}00$   & $\mathbf{+1.000.000{,}00}$ ($= c^2$) \\
4 & $4.000.000{,}0$ & $8.000.000{,}0$  & $4{,}0000$ & $1.000.000{,}00$ & $2.000.000{,}00$ & $\mathbf{+1.000.000{,}00}$ ($= c^2$) \\
\midrule
\multicolumn{7}{l}{\textbf{Trường hợp 2: Cửa sổ Hanning Praat-8 ($N+1$)}} \\
0 & $3.375.000{,}0$ & $6.750.000{,}0$ & $3{,}3750$ & $1.000.000{,}00$ & $2.000.000{,}00$ & $\mathbf{+1.000.000{,}00}$ ($= c^2$) \\
1 & $-8.931{,}5$    & $3.102.868{,}5$ & $3{,}1118$ & $-2.870{,}19$    & $997.129{,}81$   & $\mathbf{+1.000.000{,}00}$ ($= c^2$) \\
2 & $-2.431.670{,}4$ & $0{,}0$         & $2{,}4317$ & $-1.000.000{,}00$ & $0{,}00$         & $\mathbf{+1.000.000{,}00}$ ($= c^2$) \\
3 & $12.033{,}3$    & $1.602.868{,}5$ & $1{,}5908$ & $7.564{,}16$     & $1.007.564{,}16$ & $\mathbf{+1.000.000{,}00}$ ($= c^2$) \\
4 & $846.664{,}8$   & $1.693.329{,}6$ & $0{,}8467$ & $1.000.000{,}00$ & $2.000.000{,}00$ & $\mathbf{+1.000.000{,}00}$ ($= c^2$) \\
\bottomrule
\end{tabular}
\end{table}

Nhìn vào cột cuối cùng của Bảng~\ref{tab:ch05_dc_hand_calc}: ở cả hai loại hàm cửa sổ, hiệu số $r_x(\tau) - r_s(\tau)$ tại mọi độ trễ $\tau \in \{0, 1, 2, 3, 4\}$ đều bằng \textbf{chính xác tuyệt đối} $1.000.000{,}00 = c^2$. Tại $\tau = 2$ (nửa chu kỳ dao động), giá trị chuẩn hóa của tín hiệu sạch là $-1.000.000$ (tương quan nghịch cực đại, $r_s(2)/r_s(0) = -1{,}0$). Nhưng khi có DC, giá trị này bị kéo thẳng lên $0{,}0$, xóa sạch hoàn toàn dấu vết đối pha của dạng sóng!

\begin{figure}[H]
\centering
% chuong_v3/fig_ch05_dc_effect.tex
% Minh hoa tac hai cua muc DC lam lech duong tu tuong quan va pha vo nguong voicing
\begin{tikzpicture}[>=stealth, font=\small, scale=0.92, every node/.style={transform shape}]

    % Style definitions
    \tikzstyle{axisline}  = [->, thick, color=black!80]
    \tikzstyle{gridline}  = [help lines, color=gray!30, dashed]
    \tikzstyle{threshline}= [thick, dashed, color=red!80!black]
    \tikzstyle{cleanline} = [very thick, color=blue!80!black]
    \tikzstyle{dcline}    = [very thick, color=orange!90!black]
    \tikzstyle{shiftarrow}= [<->, thick, color=purple!80!black]

    % Frame boundary
    \draw[draw=blue!60!black, fill=blue!1, thick, rounded corners=6pt, dashed] 
        (-1.2, -3.2) rectangle (11.6, 5.0);
    \node[anchor=north west, font=\footnotesize\bfseries, text=blue!90!black, fill=blue!1, inner sep=2pt] 
        at (-0.9, 4.85) {Tác hại của thành phần một chiều DC lên tự tương quan chuẩn hóa};

    % Coordinate Axes
    % Origin at (0, 0)
    \draw[axisline] (0.0, 0.0) -- (10.5, 0.0) node[anchor=north west, font=\scriptsize\bfseries] {Độ trễ $\tau$ (mẫu)};
    \draw[axisline] (0.0, -2.5) -- (0.0, 4.2) node[anchor=south west, font=\scriptsize\bfseries] {Tự tương quan chuẩn hóa $r(\tau)/r(0)$};

    % Ticks on Y-axis
    \draw[thick] (-0.1, 3.6) -- (0.1, 3.6) node[anchor=east, font=\scriptsize] at (-0.15, 3.6) {$+1{,}0$};
    \draw[thick] (-0.1, 1.62) -- (0.1, 1.62) node[anchor=east, font=\scriptsize\bfseries, text=red!80!black] at (-0.15, 1.62) {$0{,}45$};
    \draw[thick] (-0.1, 0.0) -- (0.1, 0.0) node[anchor=east, font=\scriptsize] at (-0.15, 0.0) {$0{,}0$};
    \draw[thick] (-0.1, -1.8) -- (0.1, -1.8) node[anchor=east, font=\scriptsize] at (-0.15, -1.8) {$-0{,}5$};

    % Voicing Threshold Line at y = 1.62 (representing 0.45 * 3.6)
    \draw[threshline] (0.0, 1.62) -- (9.8, 1.62);
    \node[anchor=south east, font=\scriptsize\bfseries, text=red!80!black, fill=white, inner sep=1.5pt] 
        at (9.8, 1.68) {Ngưỡng hữu thanh Praat (\texttt{VOICING\_THRESHOLD} $= 0{,}45$)};

    % Grid lines
    \draw[gridline] (0.0, 3.6) -- (9.8, 3.6);
    \draw[gridline] (0.0, -1.8) -- (9.8, -1.8);

    % Ticks on X-axis (tau)
    \draw[thick] (1.5, -0.1) -- (1.5, 0.1) node[anchor=north, font=\scriptsize] at (1.5, -0.15) {$26$};
    \draw[thick] (4.5, -0.1) -- (4.5, 0.1) node[anchor=north, font=\scriptsize] at (4.5, -0.15) {$80$};
    \draw[thick] (7.5, -0.1) -- (7.5, 0.1) node[anchor=north, font=\scriptsize] at (7.5, -0.15) {$150$};
    \draw[thick] (9.5, -0.1) -- (9.5, 0.1) node[anchor=north, font=\scriptsize] at (9.5, -0.15) {$214$};

    % Shaded Region: Unvoiced region (below threshold 0.45)
    \fill[green!20, opacity=0.35] (1.5, -2.2) rectangle (9.5, 1.62);
    \node[anchor=south west, font=\scriptsize\itshape, text=green!50!black] at (1.6, -2.1) {Miền vô thanh an toàn ($r(\tau) < 0{,}45$)};

    % Curve 1: Clean Unvoiced Signal (white noise, mean=0)
    % Oscillates around 0, max peak around 0.20 (y = 0.72)
    \draw[cleanline, smooth] plot coordinates {
        (0.0, 3.6) (0.5, 1.2) (1.0, 0.1)
        (1.5, 0.5) (2.0, -0.4) (2.5, 0.65) (3.0, -0.2)
        (3.5, 0.4) (4.0, -0.55) (4.5, 0.7) (5.0, -0.1)
        (5.5, 0.3) (6.0, -0.6) (6.5, 0.55) (7.0, -0.3)
        (7.5, 0.6) (8.0, -0.2) (8.5, 0.45) (9.0, -0.4) (9.5, 0.2)
    };

    % Curve 2: Unvoiced Signal corrupted by DC offset c (+c^2 shift)
    % Lifted up uniformly by ~1.45 units (representing +c^2), pushing peaks to ~2.15 > 1.62
    \draw[dcline, smooth] plot coordinates {
        (0.0, 3.6) (0.5, 2.65) (1.0, 1.55)
        (1.5, 1.95) (2.0, 1.05) (2.5, 2.10) (3.0, 1.25)
        (3.5, 1.85) (4.0, 0.90) (4.5, 2.15) (5.0, 1.35)
        (5.5, 1.75) (6.0, 0.85) (6.5, 2.00) (7.0, 1.15)
        (7.5, 2.05) (8.0, 1.25) (8.5, 1.90) (9.0, 1.05) (9.5, 1.65)
    };

    % Vertical Shift Arrow showing +c^2
    \draw[shiftarrow] (4.5, 0.7) -- (4.5, 2.15);
    \node[anchor=west, font=\scriptsize\bfseries, text=purple!90!black, fill=white, inner sep=1pt] 
        at (4.65, 1.42) {Độ lệch $+c^2$};

    \draw[shiftarrow] (7.5, 0.6) -- (7.5, 2.05);
    \node[anchor=west, font=\scriptsize\bfseries, text=purple!90!black, fill=white, inner sep=1pt] 
        at (7.65, 1.32) {$+c^2$};

    % Critical Warning Callout
    \draw[->, thick, color=red!90!black] (2.5, 2.7) -- (2.5, 2.15);
    \node[anchor=south, font=\scriptsize\bfseries, text=red!90!black, fill=white, inner sep=1.5pt] 
        at (2.5, 2.75) {Vượt ngưỡng! Báo động giả Voiced};

    % Legend Box (Bottom)
    \node[draw=gray!60, fill=white, rounded corners=3pt, inner sep=4pt, anchor=south east] 
        at (9.6, -1.8) {
        \begin{tabular}{ll}
            \tikz\draw[cleanline] (0,0) -- (0.5,0); & \scriptsize Tín hiệu sạch (đã trừ trung bình): vô thanh đúng ($r < 0{,}45$) \\
            \tikz\draw[dcline] (0,0) -- (0.5,0);    & \scriptsize Tín hiệu có DC $+c$: bị đẩy lên $+c^2$, sai sang hữu thanh
        \end{tabular}
    };

\end{tikzpicture}

\caption{Hiện tượng mức một chiều DC $+c$ nâng bổng đường cong tự tương quan chuẩn hóa thêm một lượng $+c^2$, kéo các đỉnh nhiễu vô thanh vượt ngưỡng \texttt{VOICING\_THRESHOLD} $= 0{,}45$.}
\label{fig:ch05_dc_effect}
\end{figure}

\textbf{Đọc hình.} Hình~\ref{fig:ch05_dc_effect} so sánh trực quan hàm tự tương quan chuẩn hóa $r(\tau)/r(0)$ của một khung âm thanh vô thanh (nhiễu trắng ngẫu nhiên, $\sigma = 1000$) trước và sau khi bị nhiễm mức DC $c = +1000$. Đường màu xanh nét đậm biểu diễn tín hiệu sạch đã được trừ trung bình hoàn toàn: đồ thị dao động ngẫu nhiên quanh trục hoành $0{,}0$, đỉnh cực đại trong dải quét $\tau \in [26, 214]$ chỉ đạt $0{,}1977$, nằm hoàn toàn bên trong miền vô thanh an toàn (vùng phủ màu xanh lá, bên dưới đường đứt nét màu đỏ biểu diễn ngưỡng \texttt{VOICING\_THRESHOLD} $= 0{,}45$). Ngược lại, đường màu cam nét đậm biểu diễn tín hiệu chưa được trừ trung bình: toàn bộ quỹ đạo bị đẩy tịnh tiến lên trên một khoảng $+c^2$, khiến đỉnh cực đại bị đội lên tới $0{,}6023$. Đỉnh này vượt qua ngưỡng đỏ $0{,}45$, kích hoạt cảnh báo sai và đánh dấu sai khung vô thanh thành hữu thanh.

\textit{Tóm lại: Mức DC tĩnh $c$ sinh ra thành phần $+c^2 R_w(\tau)$ trong tự tương quan thô, và sau khi chia cho cửa sổ sẽ biến thành một độ lệch dương $+c^2$ đồng đều trên mọi độ trễ, phá hủy ranh giới nhận diện vô thanh và bắt buộc phần cứng phải có khối trừ trung bình cục bộ.}

\subsection*{5.2. Praat lấy trung bình trên đoạn nào: 426 mẫu quanh tâm khung}
\addcontentsline{toc}{subsection}{5.2. Praat lấy trung bình trên đoạn nào: 426 mẫu quanh tâm khung}

\subsubsection*{Vấn đề thực tế: Vì sao không lấy trung bình trên toàn bộ 638 mẫu cửa sổ?}
Một câu hỏi tự nhiên mà bất kỳ kỹ sư phần cứng nào cũng sẽ đặt ra khi thiết kế khối này: Khung phân tích dài $638\text{ mẫu}$, vậy tại sao không lấy tổng của toàn bộ $638\text{ mẫu}$ rồi chia cho $638$? Việc dùng một độ dài cửa sổ duy nhất cho cả phép tính trung bình lẫn phép nhân cửa sổ Hanning chẳng phải sẽ làm cho bộ đếm địa chỉ đơn giản hơn rất nhiều hay sao?

Lý do nằm ở \textbf{tính phi dừng} (non-stationarity) của tiếng nói tự nhiên. Một cửa sổ $638\text{ mẫu}$ ở tần số $16\text{ kHz}$ kéo dài tới $40\text{ ms}$. Trong khoảng thời gian $40\text{ ms}$ đó, tín hiệu tiếng nói có thể chuyển tiếp từ một phụ âm tắc vô thanh sang một nguyên âm mạnh, hoặc trải qua một biến thiên áp suất đột ngột. 

Khi áp dụng cửa sổ Hanning, trọng số ở hai mép ngoài cùng tiệm cận $0$ ($w[0] = w[637] = 1/32768 \approx 0{,}00003$), trong khi toàn bộ năng lượng hữu ích tập trung dày đặc ở vùng trung tâm ($w[318] = w[319] = 32767/32768 \approx 1{,}0$). Nếu ta tính trung bình trên toàn bộ $638\text{ mẫu}$, các đột biến biên độ ở hai đầu mép (vốn sẽ bị cửa sổ Hanning dập tắt) lại bị kéo vào và làm sai lệch giá trị mức một chiều thực tế tại vùng trung tâm. Kết quả là tại tâm khung nơi năng lượng cao nhất, mức DC cục bộ không những không được triệt tiêu mà còn bị lệch đi, để lại một thành phần dư thừa phá vỡ tính đối xứng của tự tương quan.

\subsubsection*{Giải mã mã nguồn Praat 7.0.03: Dẫn xuất khoảng lấy mẫu 426 mẫu}
Để khắc phục nhược điểm trên, Paul Boersma đã chọn một giải pháp tinh tế: \textbf{chỉ tính trung bình cục bộ trên đúng hai chu kỳ sóng dài nhất quanh tâm khung} (``one longest period to both sides''). 

Kiểm tra trực tiếp mã nguồn chuẩn của Praat phiên bản 7.0.03 (tệp \path{fon/Sound_to_Pitch.cpp}), quy trình này được định nghĩa chính xác qua các dòng mã sau:
\begin{enumerate}[leftmargin=1.5em, itemsep=2pt]
    \item \textbf{Xác định độ dài chu kỳ dài nhất (\texttt{Sound\_to\_Pitch.cpp:326}):}
    \begin{lstlisting}[language=C, caption={Độ dài chu kỳ sóng dài nhất trong Praat 7.0.03}]
nsamp_period = Melder_ifloor (1.0 / dx / pitchFloor);
    \end{lstlisting}
    Với tần số lấy mẫu $f_s = 16.000\text{ Hz}$ ($dx = 1/16000\text{ s}$) và ngưỡng đáy $F_{\text{floor}} = 75{,}0\text{ Hz}$:
    \begin{equation}
    \texttt{nsamp\_period} = \left\lfloor \frac{16000}{75} \right\rfloor = \mathbf{213\text{ mẫu}}
    \end{equation}
    \item \textbf{Định vị ranh giới mẫu bắt đầu và kết thúc (\texttt{Sound\_to\_Pitch.cpp:59--60}):}
    \begin{lstlisting}[language=C, caption={Biên giới đoạn lấy trung bình cục bộ trong Praat 7.0.03}]
startSample = rightSample - nsamp_period;
endSample = leftSample + nsamp_period;
    \end{lstlisting}
    Theo định nghĩa căn tâm của Praat (Chương 4), \texttt{leftSample} và \texttt{rightSample} là hai mẫu nằm sát hai bên điểm chính giữa của khung, thỏa mãn hệ thức \texttt{leftSample} = \texttt{rightSample} $- 1$. Do đó, số lượng mẫu nằm trong đoạn từ \texttt{startSample} tới \texttt{endSample} là:
    \begin{align}
    N_{\text{span}} &= \texttt{endSample} - \texttt{startSample} + 1 \nonumber \\
                    &= (\texttt{leftSample} + \texttt{nsamp\_period}) - (\texttt{rightSample} - \texttt{nsamp\_period}) + 1 \nonumber \\
                    &= (\texttt{rightSample} - 1 + 213) - (\texttt{rightSample} - 213) + 1 \nonumber \\
                    &= 2 \times 213 = \mathbf{426\text{ mẫu}}
    \end{align}
    \item \textbf{Tính tổng và chia trung bình (\texttt{Sound\_to\_Pitch.cpp:63--66}):}
    \begin{lstlisting}[language=C, caption={Tính tổng và chia trung bình dấu phẩy động trong Praat 7.0.03}]
for (integer i = startSample; i <= endSample; i ++)
    localMean += my z [1] [i];
localMean /= 2 * nsamp_period; // Chia cho 426 trong kieu double
    \end{lstlisting}
    \item \textbf{Trừ trung bình cục bộ trước khi nhân cửa sổ (\texttt{Sound\_to\_Pitch.cpp:72--78}):}
    \begin{lstlisting}[language=C, caption={Trừ trung bình trước khi nhân cửa sổ trong Praat 7.0.03}]
for (integer i = startSample; i <= endSample; i ++)
    thy z [1] [i - startSample + 1] = (my z [1] [i] - localMean) * window [i - startSample + 1];
    \end{lstlisting}
\end{enumerate}

\subsubsection*{Tọa độ hình học trong khung 638 mẫu và lợi thế bộ nhớ RAM vòng}
Trong cửa sổ phân tích $638\text{ mẫu}$ (chỉ số mẫu tương đối $n \in [0, 637]$), tâm khung nằm đối xứng tại vị trí giữa mẫu $318$ và mẫu $319$ (\texttt{halfnsamp\_window} $= 319$). 

Mẫu \texttt{rightSample} ứng với chỉ số tương đối $319$, còn \texttt{leftSample} ứng với chỉ số $318$. Tọa độ tương đối của đoạn $426\text{ mẫu}$ lấy trung bình cục bộ là:
\begin{align}
n_{\text{start}} &= 319 - 213 = \mathbf{106} \\
n_{\text{end}}   &= 318 + 213 = \mathbf{531}
\end{align}

Số mẫu của đoạn này là $531 - 106 + 1 = \mathbf{426\text{ mẫu}}$. Đoạn này bỏ lại đúng $106\text{ mẫu}$ ở mép trái ($0 \dots 105$) và đúng $106\text{ mẫu}$ ở mép phải ($532 \dots 637$). Hai phần mép ngoài này chiếm $106/638 \approx 16{,}6\%$ mỗi bên, chính là vùng mà trọng số cửa sổ Hanning suy giảm xuống dưới $0{,}19$ ($w[106] = w[531] \approx 6250/32768 \approx 0{,}191$).

\begin{costbox}
\textbf{Lợi ích kiến trúc phần cứng:} Ở Chương 4, ta đã thiết kế bộ đệm vòng (Ring Buffer) dung lượng $1024\text{ từ}$ có khả năng lưu trữ trọn vẹn toàn bộ $638\text{ mẫu}$ của khung hiện tại. 
\begin{itemize}[leftmargin=1.5em, itemsep=2pt]
    \item Khi bắt đầu xử lý một khung, bộ điều khiển chỉ cần phát địa chỉ đọc quét từ $\texttt{rd\_base} + 106$ tới $\texttt{rd\_base} + 531 \pmod{1024}$.
    \item Khối tính trung bình đọc trực tiếp các mẫu này từ RAM vòng mà \textbf{không cần thêm bất kỳ một khối nhớ RAM nào khác}!
    \item Chi phí bộ nhớ RAM thêm cho khối tính trung bình bằng \textbf{chính xác 0 khối M10K}.
\end{itemize}
\end{costbox}

\subsubsection*{Đối chiếu số học: Số thực dấu phẩy động của Praat so với số nguyên của FPGA}
Trong phần mềm Praat, biến \texttt{localMean} được lưu trữ và tính toán dưới dạng số thực dấu phẩy động chính xác kép (IEEE 754 double precision, 64 bit). Phép chia \texttt{localMean /= 426.0} bảo toàn phần thập phân với độ chính xác cao.

Ngược lại, trên phần cứng FPGA, nhằm tiết kiệm tài nguyên logic và duy trì thông lượng thời gian thực, thiết kế Golden v2 quy chuẩn giá trị trung bình về \textbf{số nguyên có dấu 16 bit}:
\begin{equation}
\text{mean}_{\text{FPGA}} = \left\lfloor \frac{S + 213}{426} \right\rfloor, \quad \text{với } S = \sum_{j=106}^{531} x[j]
\end{equation}

Hằng số $+213$ ($= 426/2$) đóng vai trò làm tròn số học đối xứng về số nguyên gần nhất (nearest-integer rounding). Câu hỏi đặt ra là: \emph{Việc làm tròn về số nguyên này làm lệch trung bình bao nhiêu so với số thực của Praat, và có ảnh hưởng đến kết quả F0 cuối cùng hay không?}

\textbf{1. Đánh giá cận sai số lý thuyết:}
Với mọi số nguyên tổng $S$, đặt $S = 426 \cdot q + r$, trong đó thương nguyên $q \in \mathbb{Z}$ và số dư $r \in [0, 425]$. Ta có:
\begin{itemize}[leftmargin=1.5em, itemsep=2pt]
    \item Giá trị trung bình thực dấu phẩy động của Praat là $\mu_{\text{double}} = q + \frac{r}{426}$.
    \item Giá trị trung bình số nguyên của FPGA là $\mu_{\text{int}} = \left\lfloor \frac{426q + r + 213}{426} \right\rfloor = q + \left\lfloor \frac{r + 213}{426} \right\rfloor$.
\end{itemize}
Xét hai trường hợp của số dư $r$:
\begin{itemize}[leftmargin=1.5em, itemsep=2pt]
    \item Nếu $0 \le r \le 212$: ta có $r + 213 \le 425 < 426$, do đó $\lfloor (r+213)/426 \rfloor = 0 \implies \mu_{\text{int}} = q$. Sai số tuyệt đối là:
    \begin{equation}
    |\mu_{\text{int}} - \mu_{\text{double}}| = \frac{r}{426} \le \frac{212}{426} \approx \mathbf{0{,}49765\text{ đơn vị biên độ}}
    \end{equation}
    \item Nếu $213 \le r \le 425$: ta có $426 \le r + 213 \le 638$, do đó $\lfloor (r+213)/426 \rfloor = 1 \implies \mu_{\text{int}} = q + 1$. Sai số tuyệt đối là:
    \begin{equation}
    |\mu_{\text{int}} - \mu_{\text{double}}| = 1 - \frac{r}{426} = \frac{426 - r}{426} \le \frac{426 - 213}{426} = \frac{213}{426} = \mathbf{0{,}50000\text{ đơn vị biên độ}}
    \end{equation}
\end{itemize}
Như vậy, sai số giữa trung bình số nguyên của FPGA và trung bình số thực của Praat \textbf{luôn luôn bị chặn trên bởi $0{,}5$ đơn vị biên độ} (trên thang đo $32.768$, sai số này chỉ chiếm tối đa $0{,}0015\%$).

\textbf{2. Kết quả kiểm chứng thực nghiệm trên tệp âm thanh chuẩn:}\\
Kịch bản kiểm chứng \texttt{verify\_v3/ch05\_praat\_lines.py} đã chạy phân tích đối soát từng khung trên cả hai tệp âm thanh tổng hợp mẫu:
\begin{itemize}[leftmargin=1.5em, itemsep=2pt]
    \item \textbf{Tệp âm thanh sạch (\texttt{synth\_glide\_vuv\_16k.wav}, 277 khung):} Sai số tuyệt đối cực đại đo được là $\max |\mu_{\text{int}} - \mu_{\text{double}}| = 0{,}500000$; sai số trung bình trên toàn tệp là $0{,}228954$ đơn vị; sai số phân vị $95\%$ là $0{,}465258$ đơn vị.
    \item \textbf{Tệp âm thanh có nhiễu 5 dB SNR (\texttt{synth\_glide\_vuv\_16k\_snr5.wav}, 277 khung):} Sai số tuyệt đối cực đại là $0{,}495305$; sai số trung bình là $0{,}249309$; sai số phân vị $95\%$ là $0{,}472770$ đơn vị.
\end{itemize}
Sai số cực tiểu này hoàn toàn bị triệt tiêu khi nhân với cửa sổ Hanning và tích lũy tự tương quan, bảo đảm kết quả cao độ $F_0$ cuối cùng của FPGA khớp với Praat trong phạm vi phần mười cent.

\textit{Tóm lại: Boersma chọn đoạn 426 mẫu quanh tâm khung (ứng với 2 chu kỳ của 75 Hz) để tránh hiện tượng phi dừng ở hai đầu mép cửa sổ; đoạn này tương ứng với các vị trí 106..531 trong khung 638 mẫu và được đọc trực tiếp từ RAM vòng mà không tốn thêm bộ nhớ; phép làm tròn số nguyên có sai số dưới 0,5 đơn vị biên độ.}

